% Propietario conservado: /Users/ruben/Documents/ChatGPT/jueces y controles/REVISION_ARGUMENTAL_APERTURAS_CIERRES_20260915/delivery/ARTICLE_V_ES_ARGUMENTAL_20260915/manuscrito/sections/70_radiacion_termica.tex
% RESULTADO_RECUPERADO: c40_termicas.tex del integral activo.
% Ampliación de prueba: conteo de modos, momentos térmicos, máximos
% espectrales y sustitución de la escala del calendario.
% Las hipótesis de realización se conservan explícitas.
\section{Rayonnement thermique et échelle d’information du calendrier}
\label{v:v:sec:radiacion-termica}

La statistique d’occupation fournit le lien entre les modes d’une
réalisation et leur énergie moyenne. Dans la réalisation électromagnétique,
la dispersion et la multiplicité de polarisation déterminent en outre la
densité spectrale. Leur composition permet d’obtenir les moments thermiques
et les maxima spectraux, en maintenant séparées trois opérations :
la production de l’échelle d’action, la représentation des modes et
l’évaluation de l’état thermique.


% BEGIN OWNER /Users/ruben/Documents/ChatGPT/jueces y controles/REVISION_ARGUMENTAL_APERTURAS_CIERRES_20260915/delivery/ARTICLE_V_ES_ARGUMENTAL_20260915/manuscrito/sections/69_enlace_constitutivo_velocidad.tex
% RESULTADO_RECUPERADO: funtor dimensional integral, eq:cZ-internos,
% eq:longitud-ruta y definición de ell_0; composición con artículo III.
% La prueba reunida y su control algebraico son específicos de esta contribución.
\subsection{Vitesse constitutive et réalisation cinématique de l’horloge}
\label{v:v:sec:enlace-velocidad-constitutiva}

La propagation radiative utilise la section de vitesse produite par le
transport angulaire du même état HMT. L’identification s’établit
avant le choix des unités de publication. Soient $x=A_{\mathrm{rad}}$ et
$y=C^*_{\mathrm{rad}}$ les coordonnées angulaires générées dans la
section~\ref{v:v:ant:sec:ii-angulos}, avec les canaux
\eqref{v:v:ant:eq:ii-canales}, et soit
$\mathsf R$ l’involution autoadjointe des deux feuillets, avec les projecteurs
de rang un $P_+$ et $P_-$. Sur le domaine $x>|y|$, le transport et ses
canaux sont
\begin{equation}
 T=\exp(-xI-y\mathsf R)=q_+P_++q_-P_-,
 \qquad q_+=e^{-x-y},\quad q_-=e^{-x+y},\quad 0<q_\pm<1.
 \label{v:v:eq:velocidad-transporte}
\end{equation}
La réponse des blocs temporels de longueurs $90$ et $120$ est
\begin{equation}
 \mathsf V=(I-T^{90})(I-T^{120})^{-1}
       =r_+P_++r_-P_-,\qquad
 r_\pm=\frac{1-q_\pm^{90}}{1-q_\pm^{120}}.
 \label{v:v:eq:velocidad-respuesta}
\end{equation}
La construction part des canaux hérités ; les symboles $x,y$ et
$q_\pm$ désignent leurs coordonnées et non des paramètres ajustés à une vitesse
cible.

\begin{proposition}[Identité des réalisations constitutive et cinématique]
\label{v:v:prop:identidad-velocidad}
Dans une même carte dimensionnelle, la section $c_{\mathrm{int}}$ de l’horloge
et la section $c_{\mathrm{vac}}$ de la réponse constitutive coïncident :
\begin{equation}
 c_{\mathrm{int}}=c_{\mathrm{vac}}
   =\widehat c\,u_C,
 \qquad
 \widehat c=(r_+r_-)^{-1}
            =\exp(-\mathcal E),
 \label{v:v:eq:identidad-velocidad}
\end{equation}
où
\begin{equation}
 \mathcal E=
 \log\!\bigl[(1-q_+^{90})(1-q_-^{90})\bigr]
 -\log\!\bigl[(1-q_+^{120})(1-q_-^{120})\bigr].
 \label{v:v:eq:velocidad-lector-simetrico}
\end{equation}
L’égalité conserve l’échange des feuillets.
\end{proposition}
\begin{proof}
Les facteurs qui apparaissent dans les logarithmes sont positifs. Par conséquent,
\[
 \mathcal E=\log(r_+r_-)
            =\log\det\mathsf V,
 \qquad \exp(-\mathcal E)=(\det\mathsf V)^{-1}.
\]
Le déterminant est calculé ici dans la réalisation à deux feuillets de rang
un. La définition cinématique $c_{\mathrm{int}}=\exp(-\mathcal E)u_C$
et la définition constitutive $c_{\mathrm{vac}}=(r_+r_-)^{-1}u_C$
sont, par conséquent, la même section. L’échange $P_+\leftrightarrow
P_-$ permute $r_+$ et $r_-$ et conserve leur produit. Enfin,
$\det T=e^{-2x}$ correspond au transport élémentaire, tandis que
$\det\mathsf V=r_+r_-$ correspond à sa réponse temporelle ; les deux
déterminants ont des fonctions distinctes dans cette composition.
\end{proof}

\subsubsection{Longueur de trajectoire et durée élémentaire}

Le repère dimensionnel positif $(u_S,u_C,u_T,u_Q)$ induit la base de longueur $u_L=u_Cu_T$. Pour une trajectoire enrichie $\gamma$ de $n=|\gamma|$ transitions et à canal constant, les lecteurs additifs sont
\begin{equation}
 T_{\mathrm{int}}(\gamma)=n u_T,
 \qquad L_{\mathrm{int}}(\gamma)=n\widehat c\,u_L.
 \label{v:v:eq:velocidad-longitud-reloj}
\end{equation}
Son additivité se rapporte au dénombrement et à la réalisation de chaque transition :
les chemins enrichis conservent en outre l’orientation et la mémoire.

\begin{corollary}[Normalisation de la longueur élémentaire]
\label{v:v:cor:longitud-elemental-constitutiva}
Pour $n>0$, $L_{\mathrm{int}}(\gamma)/T_{\mathrm{int}}(\gamma)
=c_{\mathrm{vac}}$. Dans la normalisation $t_0=u_T$,
\begin{equation}
 \ell_0=c_{\mathrm{int}}t_0=\widehat c\,u_L.
 \label{v:v:eq:longitud-elemental-constitutiva}
\end{equation}
Dans une carte temporelle générale $t_0=\tau_0u_T$, avec $\tau_0>0$,
l’expression correspondante est $\ell_0=\tau_0\widehat c\,u_L$.
\end{corollary}
\begin{proof}
La division de \eqref{v:v:eq:velocidad-longitud-reloj} simplifie l’entier
$n$ et utilise $u_L/u_T=u_C$. La formule élémentaire s’obtient
en multipliant la section de vitesse par $t_0$.
\end{proof}

La généralisation à un canal dépendant de l’arête conserve la formule
additive $L_{\mathrm{int}}(\gamma)=\sum_{e\in\gamma}\widehat c(e)u_L$.
Son quotient par $nu_T$ est la moyenne des vitesses de ces arêtes.
La réalisation radiative homogène emploie le canal constant déclaré
dans \eqref{v:v:eq:velocidad-longitud-reloj}.

\subsubsection{Changement de carte et normalisation adaptée}

\begin{proposition}[Covariance de l’identification]
\label{v:v:prop:covariancia-identidad-velocidad}
Soient $u_C'=d u_C$ et $u_T'=\eta u_T$, avec $d,\eta>0$.
La même vitesse et la même durée élémentaire ont pour coordonnées
\begin{equation}
 \widehat c'=d^{-1}\widehat c,
 \qquad \tau_0'=\eta^{-1}\tau_0,
 \qquad u_L'=d\eta u_L.
 \label{v:v:eq:velocidad-cambio-carta}
\end{equation}
En particulier, si $c_V=\widehat c_Vu_C^V$ et
$u_C^V=d u_C^{III}$, l’égalité des sections équivaut exactement à
$d\widehat c_V=\widehat c_{III}$.
\end{proposition}
\begin{proof}
On exprime chaque section dans les deux bases. On obtient
$\widehat c'u_C'=\widehat c u_C$ et
$\tau_0'u_T'=\tau_0u_T$. Leur produit satisfait
\[
 \tau_0'\widehat c' u_L'
 =\frac{\tau_0}{\eta}\frac{\widehat c}{d}
       d\eta u_L
 =\tau_0\widehat c u_L=\ell_0.
\]
La dernière équivalence résulte de la substitution de la relation entre les bases.
\end{proof}

Le choix adapté $d=\widehat c$ donne $u_C'=c_{\mathrm{vac}}$ et
$\widehat c'=1$. Ce choix exprime la section déjà construite dans une
base qui lui est adaptée. Le coefficient spectral dans la carte interne reste
$(r_+r_-)^{-1}$ ; l’inversion spectrale utilise cette carte interne.
Par exemple, en conservant les bases d’action et de charge, les coordonnées
constitutives se transforment comme
\[
 \widehat\mu'=d\widehat\mu,
 \qquad \widehat\varepsilon'=d\widehat\varepsilon,
 \qquad
 \widehat\mu'\widehat\varepsilon'=(\widehat c')^{-2}.
\]
Dans la carte adaptée, $\widehat\mu'=r_-/r_+$ et
$\widehat\varepsilon'=r_+/r_-$, puisque les coordonnées internes sont
$\widehat\mu=r_-^2$ et $\widehat\varepsilon=r_+^2$.

\subsubsection{Transport vers le secteur radiatif}

La dispersion de la réalisation radiative homogène s’écrit
$\omega_{\symbf{k}}=c_{\mathrm{vac}}|\symbf{k}|$.
Avec $h=2\pi\hbar$ et $\beta=(k_B\Theta)^{-1}$, sa densité spectrale
et son coefficient de Stefan--Boltzmann admettent les expressions
\begin{equation}
 u_\omega=
 \frac{\hbar\omega^3}{\pi^2(\widehat c u_C)^3}
 \frac{1}{e^{\beta\hbar\omega}-1},
 \qquad
 \sigma=\frac{\pi^2k_B^4}{60\hbar^3(\widehat c u_C)^2}.
 \label{v:v:eq:radiacion-velocidad-constitutiva}
\end{equation}
Ces substitutions transportent exactement la même section de
vitesse. Les conditions de la réalisation radiative sont conservées :
secteur bosonique libre tridimensionnel, deux polarisations transversales,
potentiel chimique nul et température positive. L’identité
\eqref{v:v:eq:identidad-velocidad} établit son antécédent constitutif ;
le dénombrement des modes et la limite thermodynamique appartiennent aux preuves
spécifiques de cette réalisation.

\paragraph{Normalisation sous amplification.}
Si le raffinement agit par $\mathsf V_m=\mathsf V\otimes I_{9^m}$,
la trace normalisée $\tau_m=\operatorname{Tr}/(2\cdot9^m)$ satisfait
\[
 \exp[-2\tau_m(\log\mathsf V_m)]
 =\exp[-\log r_+-\log r_-]=\widehat c.
\]
En effet, chaque valeur propre apparaît $9^m$ fois. La multiplicité est
simplifiée dans la trace normalisée. Le déterminant ordinaire agrandi est
$(r_+r_-)^{9^m}$ : la normalisation à deux feuillets, ou de manière équivalente la
trace normalisée, préserve le lecteur constitutif pendant ce
raffinement. Cette amplification de la représentation conserve le
transport enrichi qui en est à l’origine.

% END OWNER


\subsection{Action, calendrier et température}

L’orientation \(a\) de l’horloge est conservée et l’on abrège
\(\hbar=\hbar_a\), \(c=c_{\rm int}=c_{\rm vac}\), avec
l’identification de \eqref{v:v:eq:identidad-velocidad} dans une même carte.
L’énergie \(\varepsilon_{108}\) de cette section est
\(\epsilon_a\) de \eqref{v:v:ant:eq:ii-kb-aut-energia} ;
\(\Theta_{108}=\Theta_{\rm clk,a}\). Sa division par \(\log3\)
produit l’énergie par nat de \eqref{v:v:ant:eq:ii-kb-aut-energia-nat}.
La construction et la covariance de la paire thermique se trouvent dans
la section~\ref{v:v:ant:sec:ii-kb-aut-desarrollo}.

Nous écrivons \(h=2\pi\hbar>0\) pour la publication d’action, \(t_0>0\)
pour l’unité temporelle de sa carte dimensionnelle et \(\ell_0=ct_0\) pour
la longueur correspondante. La cellule chronologique de longueur \(108\)
détermine l’énergie
\begin{equation}
 \varepsilon_{108}=\frac{h}{108t_0},\qquad
 k_B\Theta_{108}\log 3=\varepsilon_{108},\qquad
 \beta_{108}=\frac{108t_0\log3}{h}.
 \label{v:v:eq:energia-informacion-calendario}
\end{equation}
Ici \(\beta=(k_B\Theta)^{-1}\) est le paramètre de Gibbs. L’égalité
fixe le produit énergétique de la température et du transducteur
d’entropie. L’individualisation dimensionnelle de \(k_B\) et de
\(\Theta_{108}\) utilise la section thermique et ses bases. Le logarithme
naturel mesure l’information par branche tritique ; son facteur \(\log3\) transforme
la cardinalité de la cellule en entropie sans dimension.

Nous introduisons aussi les unités associées
\begin{equation}
 t_P=\frac{108}{2\pi}t_0,\qquad
 \ell_P=ct_P,\qquad E_P=\frac{\hbar}{t_P}
 =\varepsilon_{108}.
 \label{v:v:eq:unidades-termicas-calendario}
\end{equation}
Ces égalités précisent le sens des symboles dans cette section ;
leur évaluation dans une autre carte conserve les identités entre quotients.
La longueur \(\ell_P=(54/\pi)\ell_0\) coïncide avec l’unité d’aire
de la section~\ref{v:v:ant:sec:gravity}, dans la réalisation circulaire qui y est
déclarée.

\subsection{Densité de modes et loi spectrale}

\begin{proposition}[Réalisation thermique des modes transversaux]
\label{v:v:prop:ley-espectral}
Considérons une réalisation libre dans une boîte périodique tridimensionnelle,
dont l’espace des excitations contient les vecteurs d’onde
\(\mathbf k\ne0\), avec deux modes transversaux par vecteur, une dispersion
\(\omega=c|\mathbf k|\), une énergie d’excitation \(\hbar\omega\) et
un potentiel chimique nul et \(\beta>0\). Le mode constant est maintenu
fixe comme donnée de fond. Dans la limite thermodynamique, la densité d’énergie
par fréquence angulaire est
\begin{equation}
 u_\omega(\Theta)=\frac{\hbar\omega^3}{\pi^2c^3}
 \frac{1}{e^{\beta\hbar\omega}-1},\qquad \omega>0.
 \label{v:v:eq:planck-angular}
\end{equation}
L’énergie est celle des excitations au-dessus du vide de cette
représentation.
\end{proposition}

\begin{proof}
Pour une boîte de volume \(V\), les vecteurs d’onde ont pour densité
\(V/(2\pi)^3\). Dans une couronne radiale d’épaisseur \(dk\), la densité
par volume, en incluant les deux polarisations, est
\[
 2\frac{4\pi k^2\,dk}{(2\pi)^3}=\frac{k^2}{\pi^2}\,dk
 =\frac{\omega^2}{\pi^2c^3}\,d\omega.
\]
Il s’agit de la densité limite des sommes sur le réseau périodique
sans le mode constant. Ce choix définit le produit de partition
avant la limite : chaque mode d’excitation a une énergie positive.
L’exclusion d’un seul point conserve la densité limite radiale.
L’occupation bosonique obtenue précédemment est
\(\bar n_\omega=(e^{\beta\hbar\omega}-1)^{-1}\).
Multiplier la densité par \(\hbar\omega\bar n_\omega\) démontre
\eqref{v:v:eq:planck-angular}. Près de zéro, l’intégrande énergétique est
d’ordre \(\omega^2\), et à l’infini, il décroît exponentiellement, de sorte
que son intégrale converge.
\end{proof}

La densité provient de la représentation tridimensionnelle spécifiée.
La multiplicité transversale utilisée ici a ce domaine physique ;
la dimension d’un espace harmonique sur le tore APP décrit,
séparément, la cohomologie de son complexe cellulaire.


% BEGIN OWNER /Users/ruben/Documents/ChatGPT/jueces y controles/REVISION_ARGUMENTAL_APERTURAS_CIERRES_20260915/delivery/ARTICLE_V_ES_ARGUMENTAL_20260915/manuscrito/sections/71_limite_termodinamico.tex
% FORMALIZACION_NUEVA de la justificación del límite utilizado en
% 70_radiacion_termica.tex. No introduce un generador de constantes.
\subsection{Limite thermodynamique des sommes d’occupation}
\label{v:v:subsec:limite-termodinamico}

La densité radiale précédente peut s’obtenir comme limite effective des
sommes de la réalisation périodique. L’estimation doit conserver le mode
constant fixé et contrôler uniformément tant l’origine que la queue
spectrale. Les publications d’action et de calendrier restent antérieures
à cette opération de réalisation.

\begin{lemma}[Convergence des trois observables thermiques]
\label{v:v:lem:limite-termodinamico}
Soient \(a,b>0\), \(L>0\), \(\Delta=2\pi/L\) et
\[
 F_N(r)=\frac{1}{e^{ar}-1},\qquad
 F_E(r)=\frac{br}{e^{ar}-1},\qquad
 F_Z(r)=-\log(1-e^{-ar})\quad(r>0).
\]
Pour chacune de ces fonctions, on a
\begin{equation}
 \lim_{L\to\infty}\frac{2}{L^3}
 \sum_{\nu\in\mathbb Z^3\setminus\{0\}}
 F\!\left(\frac{2\pi}{L}|\nu|\right)
 =\frac{1}{\pi^2}\int_0^\infty r^2F(r)\,dr.
 \label{v:v:eq:limite-modos-discretos}
\end{equation}
Les sommes et l’intégrale sont finies. La spécialisation
\(a=\beta\hbar c\), \(b=\hbar c\) fournit, respectivement,
les densités de nombre, d’énergie d’excitation et du logarithme de la fonction
de partition des modes transversaux.
\end{lemma}

\begin{proof}
Les trois fonctions sont continues et positives dans \((0,\infty)\).
Il existe des constantes \(C_0,C_1,d>0\), dépendant de \(a,b\) mais
indépendantes de \(\Delta\), telles que
\begin{equation}
 F(r)\leq\frac{C_0}{r}\quad(0<r\leq1),\qquad
 F(r)\leq C_1e^{-dr}\quad(r\geq1).
 \label{v:v:eq:envolventes-termicas}
\end{equation}
En effet, \(e^{ar}-1\geq ar\) contrôle \(F_N\) et \(F_E\)
près de l’origine ; en outre,
\(F_Z(r)=\log(1+F_N(r))\leq F_N(r)\leq1/(ar)\).
Pour \(r\geq1\), l’inégalité
\(-\log(1-y)\leq y/(1-y)\), avec \(y=e^{-ar}\), donne la borne
exponentielle de \(F_Z\). Les deux autres découlent de
\((e^{ar}-1)^{-1}\leq e^{-ar}/(1-e^{-a})\), en absorbant le facteur
\(r\) de \(F_E\) dans une exponentielle de taux inférieur. Ces bornes
prouvent déjà l’intégrabilité de \(r^2F(r)\).

Pour les sommes discrètes, nous regroupons les points entiers selon
\(m=\|\nu\|_\infty\geq1\). Chaque couche contient exactement
\[
 (2m+1)^3-(2m-1)^3=24m^2+2
\]
points et satisfait \(m\leq|\nu|\leq\sqrt3m\).
Fixons \(0<\delta\leq1\) et \(0<\Delta\leq1\).
Si \(\delta<\Delta\), il n’y a pas de points avec
\(0<\Delta|\nu|\leq\delta\). Sinon, en posant
\(N=\lfloor\delta/\Delta\rfloor\), la première borne de
\eqref{v:v:eq:envolventes-termicas} permet d’écrire
\begin{align*}
 \Delta^3\!\sum_{0<\Delta|\nu|\leq\delta}F(\Delta|\nu|)
 &\leq C_0\Delta^2\sum_{m=1}^{N}\frac{24m^2+2}{m}\\
 &\leq26C_0\Delta^2\sum_{m=1}^{N}m
 \leq26C_0\delta^2.
\end{align*}
Dans la deuxième ligne, on a utilisé \(m^{-1}\leq m\) puis
\(N(N+1)/2\leq N^2\) pour \(N\geq1\).
La contribution intégrale du même voisinage est aussi d’ordre
\(\delta^2\), car
\(\int_0^\delta r^2F(r)\,dr\leq C_0\delta^2/2\).

Pour \(R\geq1\), tout point avec \(\Delta|\nu|>R\) appartient
à une couche \(m>R/(\sqrt3\Delta)\). La seconde borne de
\eqref{v:v:eq:envolventes-termicas} donne
\begin{align*}
 \Delta^3\!\sum_{\Delta|\nu|>R}F(\Delta|\nu|)
 &\leq C_1e^{-dR/(2\sqrt3)}
 \Delta^3\sum_{m\geq1}(24m^2+2)e^{-d\Delta m/2}\\
 &\leq C_2e^{-dR/(2\sqrt3)},
\end{align*}
où \(C_2\) ne dépend pas de \(0<\Delta\leq1\). Pour vérifier
cette uniformité, on prend \(q=e^{-d\Delta/2}\) et on utilise
\[
 \sum_{m\geq1}q^m=\frac{q}{1-q},\qquad
 \sum_{m\geq1}m^2q^m=\frac{q(1+q)}{(1-q)^3}.
\]
La première identité provient de la série géométrique et la seconde de
l’application à deux reprises de \(q\,d/dq\), opération valable dans \(|q|<1\).
En outre,
\(1-e^{-d\Delta/2}\geq\Delta(1-e^{-d/2})\) par concavité.
Le facteur \(\Delta^3\) compense donc les dénominateurs des
deux expressions. La queue intégrale tend elle aussi vers zéro par la
borne exponentielle. Pour chaque \(\Delta>0\), ces mêmes estimations
prouvent la convergence de la somme infinie.

Dans la couronne compacte \(\delta\leq|k|\leq R\), la fonction
\(F(|k|)\) est continue et bornée. Les sommes de Riemann de maille
\(\Delta\) convergent vers son intégrale ; les deux sphères du bord
sont de mesure nulle. On fait d’abord \(\Delta\to0\), ensuite
\(\delta\to0\) et enfin \(R\to\infty\), en utilisant les bornes
uniformes précédentes. Il en résulte
\[
 \Delta^3\sum_{\nu\ne0}F(\Delta|\nu|)
 \longrightarrow \int_{\mathbb R^3}F(|k|)\,d^3k
 =4\pi\int_0^\infty r^2F(r)\,dr.
\]
Comme \(2/L^3=2\Delta^3/(2\pi)^3\), le coefficient radial final
est \(2\cdot4\pi/(2\pi)^3=1/\pi^2\), ce qui démontre
\eqref{v:v:eq:limite-modos-discretos}.
\end{proof}

Le contrôle de l’origine est effectué sur des excitations d’énergie positive.
Le mode constant reste fixé avant de former la fonction de partition :
son élimination ultérieure d’une intégrale ne remplacerait pas cette condition de
la représentation finie. La preuve établit le passage à la limite de la
réalisation déclarée et conserve la distinction entre cette réalisation et
la construction préalable de ses paramètres d’action et de calendrier.

% END OWNER


\subsection{Moments thermiques et information}

\begin{lemma}[Moment bosonique]
\label{v:v:lem:momento-bosonico}
Pour tout réel \(s>1\),
\begin{equation}
 \int_0^\infty\frac{x^{s-1}}{e^x-1}\,dx
 =\Gamma(s)\zeta(s).
 \label{v:v:eq:momento-bosonico}
\end{equation}
\end{lemma}
\begin{proof}
Pour \(x>0\), \((e^x-1)^{-1}=\sum_{m\geq1}e^{-mx}\).
Tous les termes de la somme sont non négatifs ; Tonelli permet d’intervertir la somme
et l’intégrale. Le changement \(y=mx\) donne
\(\int_0^\infty x^{s-1}e^{-mx}\,dx=m^{-s}\Gamma(s)\).
La somme converge pour \(s>1\), et l’égalité annoncée en résulte.
\end{proof}

\begin{proposition}[Densités de nombre, d’énergie et d’entropie]
\label{v:v:prop:densidades-termicas}
Dans la réalisation de la proposition~\ref{v:v:prop:ley-espectral},
\begin{align}
 n_\gamma&=\frac{2\zeta(3)}{\pi^2c^3(\beta\hbar)^3},\\
 u_\gamma&=\frac{\pi^2}{15c^3\hbar^3\beta^4},\\
 s_\gamma&=\frac{4\pi^2 k_B}{45c^3(\beta\hbar)^3}.
 \label{v:v:eq:densidades-termicas}
\end{align}
À la température du calendrier, ces relations s’écrivent
\begin{equation}
 n_\gamma\ell_P^3=\frac{2\zeta(3)}{\pi^2(\log3)^3},\quad
 \frac{u_\gamma\ell_P^3}{E_P}=\frac{\pi^2}{15(\log3)^4},\quad
 \frac{s_\gamma\ell_P^3}{k_B}=\frac{4\pi^2}{45(\log3)^3}.
 \label{v:v:eq:densidades-calendario}
\end{equation}
\end{proposition}

\begin{proof}
L’intégrale de l’occupation multipliée par la densité de modes utilise
\eqref{v:v:eq:momento-bosonico} avec \(s=3\), et l’intégrale énergétique
l’utilise avec \(s=4\). On obtient \(\Gamma(3)=2\) et
\(\Gamma(4)\zeta(4)=6\pi^4/90=\pi^4/15\).
Pour l’entropie, on part du produit de partition des modes :
\[
 \frac{\log Z}{V}=-\frac{1}{\pi^2c^3}
 \int_0^\infty\omega^2
 \log(1-e^{-\beta\hbar\omega})\,d\omega
 =\frac{\pi^2}{45c^3(\beta\hbar)^3}.
\]
La dernière égalité s’obtient en intégrant par parties après le changement
\(x=\beta\hbar\omega\). Les termes de bord s’annulent :
\(x^3\log(1-e^{-x})\) tend vers zéro aux deux extrémités.
L’identité de Gibbs \(s_\gamma/k_B=\log Z/V+\beta u_\gamma\)
prouve la troisième expression. Enfin,
\(\beta_{108}\hbar=t_P\log3\) et \(\ell_P=ct_P\) donnent les trois
normalisations du calendrier.
\end{proof}

\subsection{Flux thermique et constante de Stefan--Boltzmann}

\begin{proposition}[Flux isotrope et normalisation d’action]
\label{v:v:prop:stefan}
Le flux hémisphérique de la réalisation isotrope est
\begin{equation}
 j_*(\Theta)=\sigma\Theta^4,\qquad
 \sigma=\frac{\pi^2 k_B^4}{60\hbar^3c^2}
 =\frac{2\pi^5k_B^4}{15h^3c^2}.
 \label{v:v:eq:stefan}
\end{equation}
À la température du calendrier,
\begin{equation}
 j_*(\Theta_{108})
 =\frac{2\pi^5}{15\,108^4(\log3)^4}
 \frac{h}{\ell_0^2t_0^2}.
 \label{v:v:eq:stefan-calendario}
\end{equation}
\end{proposition}
\begin{proof}
L’énergie isotrope par unité d’angle solide est
\(u_\gamma/(4\pi)\). Le flux normal à une surface intègre
\(c\cos\theta\) sur l’hémisphère :
\[
 \frac{cu_\gamma}{4\pi}
 \int_0^{2\pi}\!d\phi\int_0^{\pi/2}\cos\theta\sin\theta\,d\theta
 =\frac{cu_\gamma}{4}.
\]
La formule pour \(u_\gamma\), \(\beta^{-1}=k_B\Theta\) et
\(h=2\pi\hbar\) démontrent \eqref{v:v:eq:stefan}.
Substituer \eqref{v:v:eq:energia-informacion-calendario} et
\(c=\ell_0/t_0\) produit \eqref{v:v:eq:stefan-calendario}.
\end{proof}

La constante \(\sigma\) est, dans cette réalisation, le coefficient
qui compose l’occupation bosonique, la densité des modes transversaux et
la projection angulaire du flux. Sa puissance quatrième provient du moment
cubique de fréquence avec la différentielle d’intégration. La lecture
du calendrier exprime ce même flux au moyen de l’action, de la longueur,
du temps et de l’information tritique.

\newpage
\subsection{Maxima spectraux et loi de déplacement}

\begin{lemma}[Racine positive des maxima spectraux]
\label{v:v:lem:raiz-wien}
Pour tout entier \(n>1\), l’équation
\(n(1-e^{-x})=x\) a une unique solution \(x_n>0\).
La fonction \(f_n(x)=x^n/(e^x-1)\) a son unique maximum positif en
\(x_n\).
\end{lemma}
\begin{proof}
Soit \(g_n(x)=n(1-e^{-x})-x\). On a
\(g_n(0)=0\), \(g_n'(x)=ne^{-x}-1\) et
\(g_n''(x)=-ne^{-x}<0\). Le maximum de \(g_n\) est atteint en
\(\log n\), où sa valeur est \(n-1-\log n>0\), tandis que
\(g_n(n)=-ne^{-n}<0\). La monotonie stricte de chaque côté du maximum
prouve l’existence et l’unicité de la racine positive. Le signe de
\(f_n'(x)\) coïncide avec celui de \(g_n(x)\) ; en outre, \(f_n\) tend
vers zéro en zéro et à l’infini. L’unicité du maximum est ainsi prouvée.
\end{proof}

\begin{proposition}[Lectures en fréquence et en longueur d’onde]
\label{v:v:prop:wien}
Les maxima de luminance énergétique par fréquence et par longueur d’onde sont
\begin{equation}
 \nu_{\max}=\frac{x_3 k_B\Theta}{h},\qquad
 \lambda_{\max}=\frac{hc}{x_5 k_B\Theta}.
 \label{v:v:eq:wien}
\end{equation}
En particulier,
\begin{equation}
 \nu_{\max}(\Theta_{108})=
 \frac{x_3}{108t_0\log3},\qquad
 \lambda_{\max}(\Theta_{108})=
 \frac{108\log3}{x_5}\ell_0.
 \label{v:v:eq:wien-calendario}
\end{equation}
\end{proposition}
\begin{proof}
La conversion \(\omega=2\pi\nu\) et le facteur angulaire isotrope donnent
\[
 B_\nu=\frac{2h\nu^3}{c^2}\frac1{e^{\beta h\nu}-1}.
\]
Avec \(\nu=c/\lambda\), la densité transformée satisfait
\(B_\lambda=B_\nu|d\nu/d\lambda|\) ; par conséquent,
\[
 B_\lambda=\frac{2hc^2}{\lambda^5}
 \frac1{e^{\beta hc/\lambda}-1}.
\]
Les changements de variable respectifs conduisent à \(f_3\) et \(f_5\).
Le lemme précédent prouve \eqref{v:v:eq:wien} ; l’échelle du calendrier
donne \eqref{v:v:eq:wien-calendario}. Le jacobien de densité explique que
les deux maxima correspondent à des racines distinctes :
\(\nu_{\max}\lambda_{\max}=c x_3/x_5\).
\end{proof}

Le déplacement de Wien est défini structurellement par le point
stationnaire de la densité spectrale choisie. Le changement de variable
transforme aussi bien l’argument que la mesure. Cette précision situe
chaque constante spectrale dans son opération génératrice et conserve sa
signification physique lors du passage entre fréquence et longueur d’onde.
